\section{Results}\label{sec:results}

This is a pre-execution study protocol; data collection has not yet begun and no quantitative results are available at the time of submission. This section reports the protocol's current status and the anticipated results that the four-month deployment is designed to produce, organised in two blocks: (a) the per-modality method-validation plan and (b) the anticipated descriptive characterisation of the multimodal corpus.

\subsection{Current status}\label{subsec:status}

The protocol has been submitted to the CEIS-UCLM (reference pending). The acquisition platform is implemented and bench-tested. The consent-and-assent campaign is to be opened upon issuance of the evaluation. Data collection is planned to start in the school term immediately following the evaluation and to run continuously for four months. The companion repository hosting the analysis code is at \url{https://github.com/marchanero/galgo-performance-protocol-maker}; the persistent identifier for the released derived-feature dataset is at \texttt{[REPOSITORY DOI --- PENDING]}.

\subsection{Method validation plan}\label{subsec:method-validation}

The Method Validation work is executed during the first two weeks of recordings and reported in a subsequent compendium article alongside the descriptive characterisation of the full four-month corpus. The protocols below specify, per modality, the measurement to be performed, the metric to be reported and the acceptance criterion against which the deployment configuration is considered fit for the descriptive characterisation that follows.

\subsubsection{Synchronisation latency analysis}\label{subsec:val-sync}

End-to-end synchronisation across the four sensing modalities is verified through a controlled LED-flash test conducted at the start of each recording week. A microcontroller whose clock is itself disciplined by the institutional NTP server drives a brief light pulse simultaneously visible to the Reolink Elite camera and physically coupled to the EmotiBit unit through a small piezoelectric actuator that triggers the accelerometer. The recorded timestamps on the camera frame, the EmotiBit accelerometer trace and the corresponding ambient-platform tick are compared against the pulse generator's reference timestamp; the difference is the per-modality offset and the dispersion across pulses is the per-modality jitter. The acceptance criterion follows the tolerances stated in \Cref{subsec:pipeline}: inter-modality alignment at or below 100~ms for the slow ambient channel and at or below 50~ms for video and EEG. The expected order of magnitude of the per-pulse offset is established in the engineering documentation of the UCLM smart-classroom platform, and the validation run is designed to confirm that figure against the specific transport-protocol configuration finally adopted \parencite{Mishra2020_mqtt,Shi2016_edge}. The validation run will report, per modality pair, the median offset and the interquartile range. The deployment is considered fit for descriptive characterisation when the upper bound of the per-pulse offset (the 75th percentile across the LED-flash test sequence) is at or below 100~ms for ambient--camera and at or below 50~ms for EEG--camera and EmotiBit--camera pairs.

\subsubsection{Physiological signal quality (EmotiBit)}\label{subsec:val-emotibit}

Signal-quality assessment for the EmotiBit layer is performed per modality and per session. For photoplethysmography-derived heart rate and heart-rate variability, the validation reports the fraction of samples retained after the project's standard quality-control rule (signal-quality index above a fixed threshold and absence of motion-flagged segments), the agreement of the EmotiBit beat-to-beat interval with a reference chest-belt electrocardiogram during a small validation block carried out under controlled laboratory conditions, and the per-session distribution of the unit's own quality channel. For electrodermal activity, the validation reports the fraction of samples surviving the low-quality flag together with the per-session count of tonic-level and phasic-event extractions retained by the decomposition pipeline. For peripheral skin temperature, the validation reports drift over the session window relative to an in-room reference temperature probe. For tri-axial accelerometry, the validation reports the per-axis dynamic range used in the session and the fraction of clipped samples. The six-unit weekly-rotation schedule and the non-dominant-forearm wear protocol are documented in \Cref{subsec:sensor-stack}; the validation block characterises whether these practical constraints leave a signal of usable quality for downstream feature extraction.

\subsubsection{EEG signal quality (Bitbrain Versatile 32, semi-dry)}\label{subsec:val-eeg}

Signal-quality assessment for the EEG layer follows the standard pipeline for semi-dry deployments. The validation will report, per session: (i) impedance distribution at instrumentation, with the per-channel acceptance threshold; (ii) the fraction of channels rejected by the offline pipeline; (iii) a broadband signal-quality estimate (frequency-domain ratio of background spectrum to identified artefact bands); and (iv) the eyeblink-artefact contamination rate before and after ICA-based artefact removal. The acceptance criterion is consistent with the dry- and semi-dry-EEG validation literature \parencite{LopezGordo2014_dry_eeg,DiFlumeri2019_dry_eeg_validation,Hinrichs2020_motion_artifacts,LauZhu2019_mobile_eeg_review}. Over four months the deployment yields approximately 16 individual EEG sessions; the protocol does not pre-commit to participant-level neurophysiological inferences and explicitly bounds EEG conclusions at the group level (\Cref{subsec:limitations}).

\subsubsection{Body-pose detection performance}\label{subsec:val-vision}

The body-pose pipeline combines a person detector and a keypoint estimator. The specific estimator (e.g., YOLOv8-Pose \parencite{Jocher2023_yolov8}, MediaPipe Holistic \parencite{Lugaresi2019_mediapipe}, OpenPose \parencite{Cao2021_openpose}, HRNet \parencite{Sun2019_hrnet}) is fixed before recordings begin and documented in the companion repository. On a held-out annotated subset of the first two weeks of recordings, the validation will report: detection mean average precision at IoU 0.5, keypoint accuracy at the standard threshold (PCK at $\alpha = 0.05$), the fraction of frames with at least the protocol-required minimum number of retained keypoints, and an explicit characterisation of detection performance under the 180\textdegree-overhead geometry and the back-row perspective \parencite{Ahuja2019_edusense}. The 72-hour deletion clock starts at the moment of feature extraction; the validation subset is processed within that window and the annotated reference frames are deleted alongside the raw footage. The retained artefact is the validation summary statistic, not the annotated frames.

\subsubsection{Environmental sensor responsivity}\label{subsec:val-ambient}

The validation will report, per channel of the UCLM ad-hoc ambient platform, the per-session drift against periodic reference-instrument checks and the datasheet accuracy of the underlying sensor (temperature and humidity from AHT21; CO\textsubscript{2}-proxy and volatile organic compound estimate from ENS160; ambient pressure from BMP280; illuminance from TEMT6000; A-weighted SPL from an I2S microphone). The validation pass also documents the 3D-printed encapsulation's passive ventilation, designed to mitigate heat-island effects inside the enclosure, and the 15-Hz nominal sampling rate. The acceptance criterion is that the per-channel drift over the deployment window stays within the datasheet's stated accuracy, complemented by the convergent validation against the cognitive-performance literature on indoor-environmental-quality channels \parencite{Wargocki2007_effects,Bako2012_classroom,Satish2012_co2,Petersen2016_classroom,Mendell2013_ventilation,Allen2016_cogfx,Twardella2012_co2_classroom}.

\subsection{Anticipated descriptive characterisation of the corpus}\label{subsec:anticipated}

Upon completion of the four-month deployment, the subsequent compendium article reports the descriptive characterisation specified in \Cref{subsec:analysis-plan}: per-modality data yield per session, per-modality missingness per window, marginal and joint distributions of the binary teacher label and the PVT change scores, and a session-level and participant-level correlation matrix (Spearman's $\rho$) between the four modalities and the two label sources. A within-teacher labelling-consistency analysis is reported via Cohen's $\kappa$ between the prompt-time annotation and a back-coded reconstruction on a spot-check subset. Group-level marginal differences between the instrumented group and the comparison group are reported as exploratory descriptors. None of these analyses constitutes a predictive claim; predictive modelling is the subject of the next compendium articles, which declare their own pre-analysis plans before any test set is touched.
